\vspace{4pt plus 4pt minus 1pt}\noindent{\kbvsecfont\fontsize{11pt}{14.5pt}\selectfont\color{kbvaccentdark} ೮.\hspace{0.45em}ಸಂಬೋಧನಾ ವಿಭಕ್ತಿ :}\hspace{0.5em}ಮತ್ತೊಬ್ಬರನ್ನು ಸಮೀಪಕ್ಕೆ ಕರೆಯುವಾಗ, ಮತ್ತೊಬ್ಬರಲ್ಲಿ ಮೊರೆಯಿಡುವಾಗ, ಮತ್ತೊಬ್ಬರೊಂದಿಗೆ ಏನನ್ನಾದರೂ ಹೇಳುವಾಗ ಈ ವಿಭಕ್ತಿ ರೂಪವನ್ನು ಬಳಸುತ್ತೇವೆ. ಸಂಬೋಧನಾ ವಿಭಕ್ತಿ ಪ್ರತ್ಯಯಗಳು ಎ, ಏ ಆಗಿವೆ.\par

\begin{kbvderiv}
\vspace{4pt plus 3pt minus 1pt}
\begin{center}\begin{adjustbox}{max width=\textwidth}
{\renewcommand{\arraystretch}{1.22}\setlength{\tabcolsep}{4pt}
\begin{tabular}{l c l}
ಉದಾ : ತಾಯೆ ! ನನ್ನನ್ನು ಹರಸು. (ತಾಯಿ + ಎ & {\color{kbvgrey}$=$} & \kbvresult{ತಾಯೆ)} \\
ತಾಯೇ ! ನನ್ನನ್ನು ರಕ್ಷಿಸು. (ತಾಯಿ + ಏ & {\color{kbvgrey}$=$} & \kbvresult{ತಾಯೇ)} \\
\end{tabular}}
\end{adjustbox}\end{center}
\vspace{4pt plus 3pt minus 1pt}
\end{kbvderiv}

\begin{kbvrule}
\noindent\kbvbullet\hspace{0.5em}ಕೆಲವು ಸಂದರ್ಭಗಳಲ್ಲಿ ನಾಮಪ್ರಕೃತಿ ರೂಪವನ್ನೇ ಸಂಬೋಧನಾ ವಿಭಕ್ತಿ ರೂಪವಾಗಿ ಬಳಸುವುದುಂಟು.
\end{kbvrule}

\noindent\hspace*{1.2em}\kbvlabel{ಉದಾ}\hspace{0.6em}ಅಕ್ಕಾ ! ಇಲ್ಲಿ ಬಾ.\par

\indent ಈ ವಾಕ್ಯದಲ್ಲಿ “ಅಕ್ಕಾ” ಶಬ್ದವನ್ನು ಸಂಬೋಧನಾ ವಿಭಕ್ತಿ ರೂಪದಲ್ಲಿ ಬಳಸಲಾಗಿದೆ. “ಅಕ್ಕಾ” ಎನ್ನುವಲ್ಲಿ ನಾಮಪ್ರಕೃತಿಯ ಕಡೆಯ ಸ್ವರ ದೀರ್ಘವಾಗಿರುವುದು ಗಮನಾರ್ಹ. (ಇಲ್ಲಿ “ಅಕ್ಕಾ” ಎನ್ನುವುದು “ಆ” ಕಾರಾಂತ, ಸ್ತ್ರೀಲಿಂಗ ಶಬ್ದ.)\par

\indent ಕೆಲವೊಮ್ಮೆ ನಾಮಪ್ರಕೃತಿಯ ಕಡೆಯ ಸ್ವರ ಹ್ರಸ್ವವಾಗಿರುವುದೂ ಉಂಟು. ಉದಾಹರಣೆಗೆ “ಮಗು, ಈ ಚೆಂಡು ತೆಗೆದುಕೊ.” (ಇಲ್ಲಿ “ಮಗು” ಎನ್ನುವುದು ಉಕಾರಾಂತ, ನಪುಂಸಕ ಲಿಂಗ ಶಬ್ದ.)\par

\begin{kbvrule}
\noindent\kbvbullet\hspace{0.5em}ಪುಲ್ಲಿಂಗ, ಸ್ತ್ರೀಲಿಂಗ, ಮತ್ತು ನಪುಂಸಕ ಲಿಂಗಗಳಲ್ಲಿನ ನಾಮಪ್ರಕೃತಿಗಳಿಗೆ ವಿಭಕ್ತಿ ಪ್ರತ್ಯಯಗಳು ಸೇರಿ ಆಗುವ ಶಬ್ದಗಳನ್ನು ಏಕವಚನ ಹಾಗೂ ಬಹುವಚನಗಳಲ್ಲಿ ಹೇಳಬಹುದು.
\end{kbvrule}

\indent ಉದಾಹರಣೆಗೆ ಅಗಸ, ರಾಣಿ, ದೇವರು — ಈ ಶಬ್ದಗಳ ವಿಭಕ್ತಿ ರೂಪಗಳನ್ನು ನೋಡೋಣ.\par

\kbvsechead{ಅಕಾರಾಂತ ಪುಲ್ಲಿಂಗ ಶಬ್ದ “ಅಗಸ” :}

\vspace{5pt plus 4pt minus 1pt}
\noindent{\renewcommand{\arraystretch}{1.24}\setlength{\tabcolsep}{4pt}
\begin{tabularx}{\textwidth}{X X X X}
{\kbvsecfont\color{kbvaccentdark}ವಿ.} & {\kbvsecfont\color{kbvaccentdark}ಪ್ರ.} & {\kbvsecfont\color{kbvaccentdark}ಏಕವಚನ} & {\kbvsecfont\color{kbvaccentdark}ಬಹುವಚನ} \\[1pt]
\arrayrulecolor{kbvline}\specialrule{0.5pt}{0pt}{2pt}\arrayrulecolor{black}
ಪ್ರ. & ವಿ. & ಅಗಸನು & ಅಗಸರು \\
ದ್ವಿ. & ವಿ. & ಅಗಸನನ್ನು & ಅಗಸರನ್ನು \\
\end{tabularx}}
\vspace{5pt plus 4pt minus 1pt}
